\subsection{Deployment protocol}

The mixed detector is exported to FP16 Core ML with non-maximum suppression
included, at confidence 0.25 and IoU 0.45 (Section~\ref{sec:ladder-method}). Mac
measurements use the same 161 images and scoring code for PyTorch and for each
Core ML processor configuration. The smartphone benchmark uses the identical
model package and three bundled 640$\times$640 images over ten repetitions. It
records median latency and p90 latency, an upper-tail value that captures slow
frames as well as typical ones. It also records load time, model outputs,
device identifier, OS, and thermal state.

The same benchmark software measures every size in the ladder on two devices: a
Mac16,8 laptop and an iPhone~17~Pro (\texttt{iPhone18,1}). Each device uses two
compute-unit settings. Compute units specify which processors Core ML may use;
the settings are CPU only and CPU with the Neural Engine. We install a Release
build of the companion application on the device and select the model
under test through a launch argument. The report is then transferred from the
device automatically, without manual transcription. The reported latency covers
inference only. It excludes camera capture, display, and application rendering,
so it is a lower bound on the cost of a live-camera pipeline.

We also measure sustained behaviour and cross-device agreement. The sustained
test runs 4{,}500 consecutive frames on the Neural Engine. We report its median,
p90, and the thermal state provided by the operating system at the end of the
run. To check cross-device agreement, we score the same images on the laptop and
the smartphone and match the two sets of boxes. This check establishes whether
the devices produce the same detections before their timings are compared.
